\clearpage
\VerdierPhysicalPageAnchor{0192}%
\begingroup\small\setlength{\parskip}{0pt}\setlength{\emergencystretch}{1.5em}
\setlength{\abovedisplayskip}{4pt plus 2pt minus 2pt}%
\setlength{\belowdisplayskip}{4pt plus 2pt minus 2pt}%

\VerdierSemanticAnchor{ch3d-2-2-5-cont}
S'il existe des entiers $n$ tels que $M\text{-}\dim(X)\leq n$, on appelle
$M$-dimension de $X$ et on note $M\text{-}\dim(X)$, le plus petit des
entiers $n$ tels que $M\text{-}\dim(X)\leq n$. S'il n'existe pas de tels
entiers, on pose $M\text{-}\dim(X)=+\infty$.

\VerdierSemanticAnchor{ch3d-2-2-6}
\medskip\noindent\textbf{Lemme 2.2.6.}\quad
\textit{Soit
$0\longrightarrow X\longrightarrow L\xrightarrow{u}Y\longrightarrow0$
une suite exacte de $\mathcal A$ telle que $L\in M$ et $X,Y\in M'$. On a :}
\[
M\text{-}\dim(X)=\sup\{M\text{-}\dim(Y)-1,0\}.
\]

\VerdierSemanticAnchor{ch3d-2-2-7}
\medskip\noindent\textbf{Proposition 2.2.7.}\quad
\textit{Soit
$0\longrightarrow X\longrightarrow Y\longrightarrow Z\longrightarrow0$
une suite exacte, $X,Y,Z\in M'$.}

\noindent\textit{a) Si
$M\text{-}\dim(Z)\leq
\inf\{M\text{-}\dim(X),M\text{-}\dim(Y)\}$, alors
$M\text{-}\dim(X)=M\text{-}\dim(Y)$.}

\noindent\textit{b) Si
$M\text{-}\dim(Y)\leq
\inf\{M\text{-}\dim(X),M\text{-}\dim(Z)\}$, alors}
\[
M\text{-}\dim(X)=
\sup\{M\text{-}\dim(Z)-1,M\text{-}\dim(Y)\}.
\]

\noindent\textit{c) Si
$M\text{-}\dim(X)\leq
\inf\{M\text{-}\dim(Y),M\text{-}\dim(Z)\}$, alors
$M\text{-}\dim(Z)=M\text{-}\dim(Y)$ ou
$M\text{-}\dim(Z)=M\text{-}\dim(Y)+1$.}

\VerdierSemanticAnchor{ch3d-2-2-7-proof}
La proposition se démontre par récurrence, en construisant un épimorphisme
de suites exactes :

\VerdierSemanticAnchor{ch3d-2-2-7-epimorphism-diagram}
\begin{center}
\begin{tikzcd}[column sep=2.8em,row sep=2.4em]
0 \arrow[r] & L_2 \arrow[r] \arrow[d] & L_1 \arrow[r] \arrow[d]
  & L_0 \arrow[r] \arrow[d] & 0 \\
0 \arrow[r] & X \arrow[r] & Y \arrow[r] & Z \arrow[r] & 0
\end{tikzcd}
\end{center}
\[
L_i\in M,\qquad 0\leq i\leq2,
\]
et en appliquant l'hypothèse de récurrence à la suite exacte noyau sur
laquelle on a un contrôle, grâce au lemme (2.2.6). Reste donc à démontrer
le lemme.

\VerdierSemanticAnchor{ch3d-2-2-6-proof}
Tout d'abord, si $M\text{-}\dim(Y)=+\infty$ (resp. $0$),
$M\text{-}\dim(X)=+\infty$ (resp. $0$). On peut donc supposer que
$M\text{-}\dim(Y)$ est fini non nul, et il s'agit de montrer que
$M\text{-}\dim(X)=M\text{-}\dim(Y)-1$. Soit :

\VerdierSemanticAnchor{ch3d-2-2-6-resolution}
\[
0\longrightarrow L_n\longrightarrow L_{n-1}\longrightarrow\cdots
\longrightarrow L_0\xrightarrow{u}Y\longrightarrow0,
\qquad n\geq1,
\]
une $M$-résolution de $Y$. Posons :

\VerdierSemanticAnchor{ch3d-2-2-6-v-complex}
\[
\begin{aligned}
V^\bullet={}&0\longrightarrow0\longrightarrow\cdots\longrightarrow0
\longrightarrow L_n\longrightarrow L_{n-1}\longrightarrow\cdots\\
&\longrightarrow L_0\xrightarrow{u}Y\longrightarrow0
\longrightarrow0\longrightarrow\cdots
\end{aligned}
\]
($L_i$ est le composant de degré $-i$),

\VerdierSemanticAnchor{ch3d-2-2-6-u0}
\[
U_0^\bullet=V^\bullet(u,1),
\qquad
p_0=p(u,1)\colon U_0^\bullet\longrightarrow V^\bullet.
\]
\endgroup
